\section{Introduction}

\todo{Week 1}

% Writing hint for Related Work paragraphs
% Sentence 1: State the main idea or research trend.
% Sentences 2-4: Cite representative studies supporting the statement.
% Sentence 5: Summarize and transition to the next topic.
% This structure ensures concise literature review and clear positioning of our work.

\paragraph{Paragraph 1, Background} 

Query information across a large volume of documents to answer a complex question (cannot answer from a single document), 3-5 related papers.

\paragraph{Paragraph 2, Problem statement} 

What is RAG, why single RAG is insufficient, and what is complex problem solving? (3-5 related papers):

\paragraph{Paragraph 3, Potential Solution} 

Explain what is multi-RAG (3-5 papers):

\paragraph{Paragraph 4, Summarize this chapter:} 
To sum up, this paper focuses on structure-aware representation learning for S1000D schema, with the goal of building a foundation for future structure-aware retrieval and generation over highly constrained technical documentation standards. This paper makes the following contributions:

\begin{itemize}
    \item We identify Text-to-S1000D XML as a challenging structured generation problem that requires not only semantic understanding but also explicit modeling of hierarchical schema constraints and domain-specific document structure.

    \item We propose a structure-aware representation learning framework for S1000D schema, using hyperbolic embedding in the Poincare Ball to capture parent-child relations, ancestor relations, hierarchical distance, and structural similarity among schema elements.
    
    \item We provide empirical evaluation through tasks such as masked tag prediction, relation classification, and distance ranking, showing that the learned representations can serve as a foundation for future retrieval, reasoning, and generation over S1000D XML.

\end{itemize}
